\section{Affine $q$-web, affine $q$-Schur, and Hall algebras}
\label{sec:affineSchur_Hall}

In this section, we show that the isomorphism between the Grothendieck groups of $\qAW \modf_\eps$ and of the module category of K-theoretic affine Hecke algebras preserves the classes of standard modules. 


\subsection{Hall algebra of linear and cyclic quivers}

Recall $\I=\Z$ or $\Z/p\Z$ from \eqref{eq:I}. We denote by $\peni$ and $\penp$ the linear and cyclic quiver with arrow sets $\Omega =\{(i,i+1)\mid i\in \I\}$.
%We drop the index to write $\pen$ to refer to both if we do not need to specify them. 
Sometimes we allow $p=\infty$ to treat both cases simultaneously. 

Let $\bbF_{\qq^2}$ be the finite field of $\qq^2$ elements. For an $\I$-graded $\bbF_{\qq^2}$-vector space $V=\oplus_{i\in\I} V_i$, denote by $E_V=\oplus_{(i,j)\in\Omega} \Hom(V_i, V_j)$ the space of nilpotent representations of $\pen_p$. The group $G_V =\prod_{i\in \I} GL(V_i)$ acts on $E_V$ by conjugation. Denote by $S_i$ the 1-dimensional simple module supported at $i\in\I$ with dimension vector denoted by $\epsilon_i$. For each pair $(i,\ell)\in \I\times \N$ there exists a unique indecomposable $\pen_p$-module $S_{i,\ell}$ of length $\ell$ and head $S_i$. In this way, each $\pen_p$-module $M$ is isomorphic to a unique
\[
M \cong \bigoplus_{i,\ell} m_{i,\ell} S_{i,\ell}.
\]
Denote by $[i;\ell)$, for $i\in\I, \ell \in\Z$, the segment $[i,i+\ell-1]$ in case $\I=\Z$ and cyclic segment equal to the modulo $p$ projection of $[i_0,i_0+\ell-1]$ for any preimage $i_0\in \Z$ of $i$ in case $\I=\Z/p\Z$. Then the isoclasses of $\pen$-modules are parametrized by the set of multi-segments
\begin{align}
    \label{cMcS}
    \cM\cS =\Big\{\bfm :=\sum_{i\in\I,\ell\in\Z} m_{i,\ell} [i;\ell) \mid m_{i,\ell} \in \N, \text{all but finitely many are }0 \Big\}.
\end{align}
Denote by $|\bfm|=\sum_{i\in\I,\ell\in\Z} \ell m_{i,\ell}$, $\cM\cS_n=\{\bfm\in \cM\cS\mid |\bfm|=n\}$, and so $\cM\cS=\cup_n\cM\cS_n$.

For an $\I$-graded $\bbF_{\qq^2}$-vector space $V=\oplus_{i\in\I} V_i$, let $\C_G(E_V)$ be the set of $\C$-valued $G_V$-equivariant functions on $E_V$. Fix an $\I$-graded $\bbF_{\qq^2}$-vector space $V_{\bfd}$ of dimension $\bfd \in \N^\I$. Given $\bfd', \bfd'' \in \N^\I$ such that $\bfd'+\bfd''=\bfd$, we have the following diagram
\[
E_{V_\bfd'} \times E_{V_\bfd''} \stackrel{p_1}{\longleftarrow} E \stackrel{p_2}{\longrightarrow} F \stackrel{p_3}{\longrightarrow} E_{V_{\bfd}}.
\]
Here $E$ denotes the set of triples $(x,\phi,\psi)$ with $x\in E_{V_\bfd}$ and
\[
0\longrightarrow V_{\bfd'} \stackrel{\phi}{\longrightarrow} V_{\bfd} \stackrel{\psi}{\longrightarrow} V_{\bfd''} \longrightarrow 0
\]
an exact sequence of $\I$-graded vector spaces such that $\phi(V_{\bfd'})$ is stable by $x$, and $F$ denotes the set of pairs $(x,U)$, where $U\subset V_\bfd$ is an $\I$-graded $x$-stable subspace of dimension $\bfd'$. 

The Hall algebra $\bold{Hall}_{p,\bbF_{\qq^2}} :=\oplus_{\bfd} \C_G(E_{V_\bfd})$ is endowed with the following multiplication \cite{Lus91, Lus98}: for $f \in \C_G(E_{V'_\bfd}), g\in \C_G(E_{V''_\bfd})$,
\begin{align}
    \label{product}
    f \circ g =\qq^{-m(\bfd'',\bfd')} (p_3)_! h \in \C_G(E_{V_\bfd}),
\end{align}
where $h \in \C(F)$ is the unique function such that $p_2^*(h) =p_1^*(f,g)$ and $m(\bfd'',\bfd')= \sum_{i\in\I} d_i'd_i'' +\sum_{i\in \I} d_i''d_{i+1}'$. 

For $\bfm \in \cM\cS$ with $\underline{\dim} (\bfm) =\bfd$, we denote by $\cO_\bfm$ the $G_{V_\bfd}$-orbit in $E_{V_\bfd}$ consisting of representations in the isoclass $\bfm$. Let $f_\bfm$ be $\qq^{-\dim \cO_\bfm}$ times the characteristic function of $\cO_\bfm$.

It is known (due to Jin Yun Guo) that the structure constants for the Hall algebras $\bold{Hall}_{p,\bbF_{\qq^2}}$ with respect to the basis $\{f_\bfm\}$ are polynomials in $\qq$. This allows one to define the generic Hall algebra $\Hall$ over $\cA=\Z[v,v^{-1}]$ for an indeterminate $v$ such that $\Hall|_{v\mapsto \qq^{-1}} =\bold{Hall}_{p,\bbF_{\qq^2}}$. We shall use the same notation $f_\bfm$ to denote the corresponding elements in $\Hall$. For $\bfm \in \cM\cS$, we let
\begin{align*}
    b_\bfm =\sum_{i, \bfn} v^{-i+\dim \cO_\bfm -\dim \cO_\bfn} \dim \cH^i_{\cO_\bfn} (\text{IC}_\bfm)f_\bfn,
\end{align*}
where $\cH^i_{\cO_\bfn} (\text{IC}_\bfm)$ is the stalk over a point of $\cO_\bfn$ of the $i$th IC sheaf of the closure $\overline{\cO_\bfm}$ of $\cO_\bfm$. Then $\bfB :=\{b_\bfm\mid \bfm \in \cM\}$ is the canonical basis of $\Hall$. There is a unique bar involution on $\Hall$ such that $\overline{v}=v^{-1}$, and $\overline{b_\bfm}=b_\bfm$, for all $\bfm$; cf. \cite{Lus98, VV99}. 

\begin{proposition}  \cite{VV99, DF15} 
\label{prop:generators}
    The Hall algebra $\Hall$ is generated by $f_\bfi$, for all $\bfi\in \Z\I$. 
\end{proposition}

\subsection{Quantum group $\U_v(\widehat\gl_p)$}

Denote by $\U(\widehat\sll_p)$ the quantum group of type $A_{p-1}^{(1)}$, i.e., the $\C(v)$-algebra generated by $E_i, F_i, K_i^{\pm 1}$, for $i\in \I=\Z/p\Z$. Let $\UA(\widehat\sll_p)$ be the $\cA$-subalgebra generated by the divided powers $E_i^{(r)}, F_i^{(r)}, K_i^{\pm 1}$, for $i\in \I=\Z/p\Z$ and $r\ge 1$.
There is an embedding $\UA(\widehat\sll_p) \rightarrow \Hall$ such that $F_i^{(r)} \rightarrow f_{r[i]}$, for all admissible $i, r$. 

We shall need to enlarge the $\U_v(\widehat\sll_p)$ to the $\C(v)$-algebra 
\begin{align} \label{eq:affine_glp}
\U_v(\widehat\gl_p):=\U_v(\widehat\sll_p) \cdot \Heis,
\end{align}
where the Heisenberg algebra $\Heis$, which  commutes with $\U_v(\widehat\sll_p)$, is the $\C(v)$-algebra generated by $B_k$, for $k\in \Z^\times$, subject to the relations \cite{KMS95, LT96, Ugl00} 
\[
[B_k, B_m]= k\delta_{k,-m} \frac{[kp]_v[k\ell]_v}{[k]_v^2},
%\frac{(1-v^{-2kp})(1-v^{-2k\ell})}{(1-v^{-2k})^2}  \text{id}.
\]
We write
$\Heis =\Heis^- \otimes \Heis^+$, where $\Heis^+$ and $\Heis^-$ are commutative subalgebras of $\Heis$ generated by $B_r$ for $\pm r>0$, respectively. Denote $\U^-_v(\widehat\gl_p) = \U^-_v(\widehat\sll_p) \cdot\Heis^-$.

There is a $\C(v)$-algebra isomorphism (cf. \cite{X97, Sch00, DF15})
\begin{align}
    \label{eq:isom}
    \xi: \C(v) \otimes_\cA\Hall \stackrel{\cong}{\longrightarrow}  \U_v^-(\widehat\gl_p).
\end{align}
We then define an $\cA$-form $\U_\cA^-(\widehat\gl_p)$ as the image of $\Hall$ under the isomorphism $\xi$:
\[
\xi: \Hall \stackrel{\cong}{\longrightarrow} \U_\cA^-(\widehat\gl_p), 
\]
and it inherits a canonical basis from $\Hall$ this way. We have $\U_\cA^-(\widehat\gl_p)\supset \U_\cA^-(\widehat\sll_p)$ with compatible canonical bases since Lusztig \cite{Lus90, Lus98} defined the canonical basis of $\U_\cA^-(\widehat\sll_p)$ geometrically by identifying it with the composition subalgebra of $\Hall$.


\subsection{Connecting $\Hall$ to $\K^*$}

Recall the geometric (K-theoretic) construction of the affine $q$-Schur algebras from \eqref{KSchur}, which we shall set $N=n$ from now on. We review some well known constructions in geometric representation theory; cf. \cite{CG97, GV93, Lus91, Lus98, AJL}. The simple modules of the convolution algebra $(\KS(N,n), \star)$ are parametrized by $GL(n)$-orbits of pairs $(s,x) \in GL(n) \times \gl_n$, where $s$ is semisimple, $x$ is nilpotent, and $sxs^{-1} =z^{-2}x$. The $GL(n)$-orbits of the nilpotent elements $x$ are parametrized by the Jordan forms or partitions $\lambda \in \Par(n)$. The $GL(n)$-orbits of pairs $(s,x)$, where the eigenvalues of $s$ are assumed to be in $z^{2\Z}$, are in bijection with the isoclasses of the $n$-dimensional nilpotent representations of $\peni$ if $z$ is generic and of $\penp$ if $z =\varepsilon$ (where $\varepsilon^2$ is a primitive $p$th root of $1$); hence the orbits $O_\bfm$ are naturally parameterized by multi-segments $\bfm\in\cM\cS$;  cf. \eqref{cMcS}.  

Associated to each such orbit $O_\bfm$, there exists a standard module $M_\bfm$ of the convolution algebra $(\KS(n,n), \star)$ defined as Borel-Moore homology of the fixed point locus of a Springer fiber $H_*(\cF_x^s)$ \cite{CG97, GV93}. Denote by $[\KS(n,n)\modf]$ the Grothendieck group of finite-dimensional $\KS(n,n)$-modules over $\C$. Then $[\KS(n,n)\modf]$ admits two bases: 
\[
\{[M_\bfm] \mid \bfm\in \cM\cS_n\},
\quad \text{ and } \quad
\{[L_\bfm] \mid \bfm\in \cM\cS_n\}, 
\]
with uni-triangular transition matrices between them; here $L_\bfm$ denotes the simple modules. Denote by $[\KS(n,n)\modf]^*$ the restricted dual, which is the $\Z$-span of the dual basis $[L_\bfm]^*$, for $\bfm\in\cM\cS$, where $[L_\bfm]^* ([L_{\bfm'}]) =\delta_{\bfm,\bfm'}$. Similarly, we denote by $\{[M_\bfm]^*\}$ the $\Z$-basis dual to $\{[M_{\bfm}]\}$. We denote 
\begin{align} \label{eq:KK}
{\K} :=\bigoplus_{n=0}^\infty [\KS(n,n)\modf],
\qquad
{\K}^* :=\bigoplus_{n=0}^\infty [\KS(n,n)\modf]^*.
\end{align}

\begin{theorem} [\text{\cite[Theorems 6.6, 8.4]{GV93}, \cite[\S12.3]{VV99}}]
\label{thm:Hall=Schursimple}
Let $p\le \infty$. 
    The $\Z$-linear isomorphism 
    \begin{align*}
        \imath_p: \HallZ \longrightarrow {\K}^*,  \qquad
        f_{\bfm} \mapsto [M_{\bfm}]^*
        \textrm{ for } \bfm \in \cM\cS,
    \end{align*}
sends the canonical basis to the basis dual to the simple modules, i.e., $\imath_p(b_\bfm) =[L_{\bfm}]^*.$
\end{theorem}


\subsection{Matching the standard modules}

We now bring back the affine $q$-webs into the consideration. 

\begin{proposition}
\label{prop:Morita}
    For any $N\ge n$, the algebras $\KS(N,n), \ASch(N,n)$ and $\qAW(n)$ are all Morita equivalent. (We only need the case when $N=n$.)
\end{proposition}

\begin{proof}
Thanks to $N\ge n$, each summand in $\hatT_{n,n}$ appears as a summand in $\hatT_{N,n}$; cf. \eqref{T:module}. Hence by the definition \eqref{def:affineSchur}, the affine $q$-Schur algebra $\ASch(N,n)$ is Morita equivalent to $\ASch(n,n)$. It follows by Proposition \ref{prop:SameS} that $\KS(N,n)\cong \ASch(N,n)$ is Morita equivalent to $\KS(n,n) \cong \ASch(n,n)$.

On the other hand, the algebra $\ASch(n,n)$ is Morita equivalent to $\qAW(n)$ by \cite[Proposition 4.18]{SSW25}. We are done. 
\end{proof}

\begin{rem}
\label{rem:composition}
    One can define a variant $'\!\qAW$ of $\qAW$ whose objects are compositions instead of strict compositions; cf. \cite{Bru25} in the setting of the $q$-web category. Then the path algebra for $'\!\qAW(n)$ is isomorphic to $\ASch(n,n)$, and hence isomorphic to $\KS(n,n)$. 
\end{rem}

\begin{definition} 
A finite dimensional $\qAW$-module $M$ is called $\eps$-integral  if it factors through the quotient $\qAW\rightarrow\qW_\bfu$ with $\bfu=(\eps^{2s_1}, \ldots, \eps^{2s_\ell})$, for some $s_i\in \I$. Let $\qAW\modf_\eps$ denote the category of $\eps$-integral $\qAW$-modules.  Categories $\hwS(n)\modf_\eps$, $\qAH\modf_\eps$, and $\AHe(n)\modf_\eps$ are defined similarly.
\end{definition}

The study of general finite-dimensional $\hwS(n)$-modules can be reduced through Morita equivalences to 
$\hwS(m)\modf_\eps$, for various $\eps$ and $m$. 

For any $\bfm=\sum_{i\in \I,\ell\in \Z}m_{i,\ell} [i;\ell) $, let $\lambda_{\bfm}$ be the multipartition such that each part is a column partition and for each $i$ there are $m_{i,\ell}$ times of  $(1^{\ell})$ for $i\in \I,\ell\in \Z $. Here we fix some order of the component of $\lambda_\bfm$. 
Set $\bfu^\diamond=((\eps^{2(i-\ell+1)})^{m_{i,\ell}})_{i\in \I,\ell\in \Z}$ such that the order corresponds to the order of component of $\lambda_\bfm$. Denoting $n=|\lambda_{\bfm}| =|\bfm|$, we define the {\em standard modules} of $\hwS(n)$:
 \begin{align}  \label{eq:std}
\Delta_\bfm:= \pi_{\bfu^\diamond}(\Delta^{\lambda_\bfm}).
\end{align}
These are always induced modules by Theorem \ref{thm:induced}. 
Similarly, the $\AHe(n)$-modules  
\[
\texttt S_\bfm:= \pi_{\bfu^\diamond}( S^{\lambda_\bfm})
\]
are induced modules. 

The Morita equivalence in Proposition \ref{prop:Morita} (and Remark \ref{rem:composition}) induces a canonical isomorphism 
\begin{equation}
\label{equ:Moritiso}
    \natural: [\qAW \modf_\eps]\longrightarrow  \K,
\end{equation}
which matches the simples. 
Similarly, we have an isomorphism $\natural: [\qAH\modf_\eps] \rightarrow\bigoplus_{n=0}^\infty[\KH(n)\modf]$. The following is built on the induction theorem from \cite{KL87} and \cite{Ar96}.

\begin{theorem} \label{thm:ind=std}
Let $p\le \infty$. The isomorphism $\natural$ in \eqref{equ:Moritiso} matches the standard modules, i.e., 
$\natural([\Delta_\bfm] )= [M_\bfm]$, for any $\bfm\in\cM\cS$. ($\natural$ matches the simples as well.)
\end{theorem}

\begin{proof}
First assume that $p=\infty$. In this case we have a commutative diagram 
\begin{equation*}  
\begin{tikzcd}
{[}\hwS(n)\modf_\eps{]}\ar[r,"\natural"]\ar[d,"1_{(1^n)}"] & {[}\KS(n,n)\modf{]} \ar[d,"1_{(1^n)}"]
\\
{[}\AHe(n)\modf_\eps{]} \ar[r,"\natural"] & {[}\KH(n)\modf{]}
\end{tikzcd} 
\end{equation*}
where the Schur functors $1_{(1^n)}$ are  isomorphisms. By Corollary \ref{cor:idemonindu}, the left-column Schur functor satisfies that
$1_{(1^n)} ([\Delta_\bfm]) = [\texttt S_\bfm]$, while the right-column Schur functor satisfies that $1_{(1^n)} ([M_\bfm]) = [N_\bfm]$; cf. \cite[Proposition 18]{V98}. Here $N_\bfm$ is the standard $\KH(n)$-module defined as in \cite[\S 6.2]{AJL}, which is related to the standard module used in \cite{Ar96} via an involution; also see Remarks \ref{rem:sign}--\ref{rem:opposite}.  Then we have $\natural([\Delta_\bfm] )= [M_\bfm]$ since $\natural([ {\texttt S}_\bfm] )= [N_\bfm]$ by Proposition \ref{prop:induhe} and \cite[Theorem 3.2]{Ar96} (also see \cite[Theorem~ 6.2]{KL87}).
Finally, the identity in case for $p>0$ follows from that for $p=\infty$ since $\Delta_\bfm$ and $M_\bfm$ are defined over any field and the equality in the Grothdendick group holds in specializations of $q\mapsto \eps$.
 \end{proof}
 
 \begin{rem}
 \label{rem:donotorder}
 It follows from Theorem \ref{thm:ind=std} that $[\Delta_\bfm]$ does not depend on the order of components in $\lambda_\bfm$ in the Grothendieck group since $[M_\bfm]$ does not. 
 \end{rem}
 
\begin{rem} \label{rem:sign}
The induced modules ${\texttt S}_\bfm$ here are induced from the sign representation of the parabolic subalgebra, while the induced modules used in \cite{Ar96} is from the trivial representation. Let ${\texttt S}'_\bfm := \pi_{\bfu'}(  S^{\lambda_\bfm'})$, where  $\bfu'=((q^{2i})^{m_{i,\ell}})_{i\in \I,\ell\in \Z}$ and $\lambda_\bfm'$ is the multipartition such that each component is a row partition and for each $i$ there are $m_{i,\ell}$ times of  $(\ell)$ for $i\in \I,\ell\in \Z$.
 Here, the order of the components of $\lambda_\bfm'$ is chosen to be compatible with that of $\bfu'$. 
 Then by 
\eqref{equ:antiind}, we have 
\begin{equation}
\label{equ:signtriv}
   {\texttt S}_\bfm^\tau\cong  {\texttt S}'_{\bfm^\tau}. 
\end{equation}
where $\bfm^\tau=\sum_{i\in \I, \ell\in \N}m_{i,\ell} (\ell, -i]$.
See also \cite[\S 6.2]{AJL}.
Note that  by \eqref{equ:shuffle}, the isoclasses of $\texttt S'_{\bfm^\tau}$ and $\texttt S_\bfm$ do not depend on the ordering of the tuple $\lambda_\bfm$ and $\lambda_{\bfm^\tau}'$. Then $\texttt [S_{\bfm }]$ also corresponds to $f_\bfm$   since $\tau$ corresponds to the anti-automorphism  $\rho$ of the Hall algebra in \cite[\S 2.3]{AJL} and $\rho (f_\bfm)=f_{\bfm^\tau}$. 
\end{rem}

\begin{rem}
\label{rem:opposite}
The discrepancy between the induced representations from the sign and trivial representations noted in Remark~\ref{rem:sign} has its root in the choice of opposite orientations on the cyclic quiver. In this paper we use the clockwise orientation ($i\mapsto i+1$) following \cite{VV99, DF15}. If we choose to work with the Hall algebra with opposite orientation as in \cite{Ar96, LTV99} and define $f_\bfi$ via  {\em right} multiplication (with multiplication formula derived from \cite{DF15, FL19} accordingly), then we expect the new formulation will match with the induced modules from the trivial representation. 
\end{rem}

\subsection{Affine $q$-web categorifies the Hall algebra}

The $q$-web category has a great advantage over $\K$ as restriction and induction functors are defined. We define linear operators 
\begin{align}
    \label{euq:defoffi}
    \begin{split}
    f_\bfi: [\qAW \modf_\eps]^* &\longrightarrow [\qAW \modf_\eps]^*, \\
    \phi \mapsto f_\bfi \phi &\qquad \text{ such that }
    f_\bfi \phi ([M]):= \phi ([\bfi\text{-Res}(M)]).
    \end{split}
\end{align}
Let $D_\bfm$ be the irreducible module in $[\qAW \modf_\eps]$ such that  
$\natural ([D_\bfm])= L_\bfm$, where $\natural$ is the isomorphism in \eqref{equ:Moritiso}. Denote $\HallZ =\Z\otimes_\cA \Hall$, the specialization at $v=1$. (Similar notations with subscript $\Z$, e.g., $\Uglzhalf$, will be freely used for the specialization at $v=1$ of $\cA$-algebras and $\cA$-modules.) 

The identifications in Proposition \ref{prop:Morita} and Theorem \ref{thm:ind=std} allow us to formulate an improved version of Theorem \ref {thm:Hall=Schursimple} in terms of the affine $q$-web category. It follows from Theorem \ref{thm:ind=std} that  
$\{[\Delta_{\bfm}]~|~ \bfm\in \cM\cS\}$ is a basis of $[\qAW \modf_\eps]$. Denote by $\{[\Delta_{\bfm}]^*\}$ the basis in $[\qAW \modf_\eps]^*$ dual to the basis $\{[\Delta_{\bfm}]\}$. 

 \begin{theorem}
 \label{thm:Cat_affine}
 Let $p\le \infty$.
 \begin{enumerate}
     \item 
     There exists a $\Uglzhalf$-module isomorphism
\begin{align}
\label{eq:simple=dCB:p}
\jmath_p: \HallZ &\longrightarrow [\qAW \modf_\eps]^*,
\qquad
f_\bfm \mapsto [\Delta_{\bfm}]^*,
\end{align}
where $f_\bfi$ acts on $\Hall$ by the multiplication with the corresponding semisimple generator and 
acts on $[\qAW \modf_\eps]^*$ by \eqref{euq:defoffi}.
\item The isomorphism $\jmath_p$ sends the canonical basis to the dual basis of the simple modules, i.e., $\jmath_p(b_\bfm)=[D_\bfm]^*$.
 \end{enumerate}
\end{theorem}
 
\begin{proof}
(1) Clearly $\jmath_p$ defined by \eqref{eq:simple=dCB:p} is a $\Z$-linear isomorphism as it matches bases. Then it remains to show that 
    \begin{equation}
    \label{equ:isomodule}
        \jmath_p(f_{\mathbf i}f_\bfm) = f_{\mathbf i} ([\Delta_{\bfm}]^*)
    \end{equation} 
for all semisimple generators $ f_{\mathbf i}$ and all multi-segments $\bfm$.

We assume $p=\infty$ first.
Write $f_{\mathbf i}f_\bfm=\sum c_{\bfm',\bfm}f_\bfm'$
and $f_{\mathbf i} [\Delta_{\bfm}]^* =\sum d_{\bfm',\bfm}  [\Delta_{\bfm'}]^*$.
Write a multi-segment as $\bfm_A=\sum_{i<j\in \Z}a_{i,j}[i,j)$ for some  upper integral  triangular  matrix $A=(a_{ij})$. Let $\mathbf T_{\bfi}$
be the set of all upper triangular integral matrices $T$ such that $\sum_{j}t_{kj}=i_k$ for all $k$. Then by \cite[Proposition 4.2]{DF15}
we have 
\[
f_\bfi f_{\bfm_A}= \sum_{ T\in \mathbf T_\bfi} \prod_{i,j\in\Z} \binom{a_{ij}+t_{ij}-t_{i-1,j}}{t_{ij}} f_{\bfm_{A+T-\tilde T^+} }   
\]
where 
$ \tilde T$ is obtained by descending each row of $T$  to the next row and $\tilde T^+$ is the upper triangular part of $\tilde T$. 


On the other hand, $d_{\bfm',\bfm} $ is determined by the $\mathbf i$-Res such that 
\begin{equation}
\label{equ:i-RES}
  [\iRes(\Delta_{\bfm'})]= \sum d_{\bfm',\bfm}[\Delta_{\bfm}].  
\end{equation}
Using Proposition \ref{prop:IndRes} and comparing the coefficients (where induced modules in the Grothdendick group for different orderings are identified, see Remark \ref{rem:donotorder}) we  see that 
\[d_{\bfm_{A+T-\tilde T^+},\bfm_A}= \prod_{i,j\in\Z} \binom{a_{ij}+t_{ij}-t_{i-1,j}}{t_{ij}}.\]
For example, if $\bfi=\lambda_i\alpha_i$, then by Proposition \ref{prop:IndRes} and Remark \ref{rem:donotorder}, 
any $\bfm'$ appearing in \eqref{equ:i-RES} is obtained from $\bfm=\sum_{i,j}a_{i,j}[i,j)$ by deleting $\lambda_i$ times of $i$ from the segments starting with $i$, i.e., 
$\bfm'=m_{A+T-\tilde T^+}$ for some $T\in \mathbf T_\bfi$, if $m=m_A$.
Furthermore, we have 
\[d_{\bfm',\bfm}=\prod_{j}\binom{a'_{ij}}{t_{ij}} =\prod_{j\in \Z} \binom{a_{ij}+t_{ij}-t_{i-1,j}}{t_{ij}}=c_{\bfm',\bfm}.\]
For general $\bfi=\sum_{i}\lambda_i\alpha_i$, the equality $d_{\bfm',\bfm}=c_{\bfm',\bfm}$ follows from the case $\bfi=\lambda_i\alpha_i$ and the induction since 
$f_\bfi=\prod_{i}f_{\lambda_i\alpha_i}$ by \cite[Remark 6.1]{VV99}, where the order of the product is from $\infty$ to $-\infty$. This proves \eqref{equ:isomodule} for  $p=\infty$. 

Below we shall prove the result for $p>0$ by using the result proved above for $p=\infty$. 

% It remains to establish the claim that $\{[\Delta_{\bfm}]\mid \bfm\in \cM\cS\} $ is a basis of $[\qAW \modf_\eps]$.

% We first prove the claim when $p=\infty$.
% In this case, we have an isomorphism between 
% \[
% 1_{1^m}: [\qAW(m)\text{-mod} ] \longrightarrow [\AHe_m\text{-mod}],
% \qquad
% [M]\mapsto [1_{1^m}M].
% \]
% since the classifications of simple modules for both algebras are all indexed by the set of all  multi-segments (of total length $m$). 

% On the other hand, it is known that each standard module $\tilde S'_\bfm $ has a simple head and all other composition factors are smaller under some partial order \cite{LTV99}. This implies 
% $ \{ [\tilde S'_\bfm]\mid \bfm\in \cM\cS\} $ is a basis of $[\AHe_m\text{-mod}]$ and hence $ \{ [\tilde S_\bfm]\mid \bfm\in \cM\cS\} $ is  also a basis of $[\AHe_m\text{-mod}]$ by \eqref{equ:signtriv}. 
% Then the claim follows since 
% $ 1_{1^m} S_\bfm= \tilde S_\bfm$ by Corollary \ref{cor:idemonindu}.

For $p>0$, we let $\bfW(n)$ be the $\Z$-submodule of 
$ [\qAW(n)\text{-mod}_\eps]$ spanned by $\{[\Delta_\bfm] \mid \bfm\in \cM\cS_n\}$. Let $\bfW=\oplus_{n\in \N}\bfW(n)$ and let $\bfW^*=\oplus_{n\in \N}\Hom(\bfW(n),\C)$ be the graded dual.
For $p=\infty$, we add a subscript $\infty$ to denote the counterparts by $\bfW_\infty$ and $ \bfW_\infty^*$. 

We have shown above that $\bfW_\infty =[\qAW \modf_\eps]$ and that $\{[\Delta_\bfm]\mid \bfm\in\cM\cS_\infty\}$ is a basis of $\bfW_\infty$.
Clearly we have a surjective $\Z$-linear  map 
\[
\varpi: \bfW_\infty \longrightarrow \bfW,
\qquad
\varpi([\Delta_\bfm])=[\Delta_{\bar\bfm}], 
\]
where $\bar\bfm$ is the multi-segment obtaned from $\bfm$ modulo $p$. Dualizing this gives us an injective $\Z$-linear map 
\begin{align*}
  \varpi_*:  \bfW^* &\longrightarrow \bfW_\infty^*, 
  \qquad
    \phi \mapsto \phi \circ \varpi. 
\end{align*}

%Recall the action of (the generators $f_\bfi$ of) $\U_{p,\Z}^-$ on $\bfW^*$ is given by $\phi \mapsto f_\bfi \phi$ with $f_\bfi \phi ([M])= \phi ([\bfi\text{-Res}(M)])$
 Thanks to Proposition \ref{prop:IndRes} and Theorem \ref{thm:induced}, we have $\varpi_*(f_\bfi \phi) ([\Delta_\bfm]) =\sum _{\bfj:\bar\bfj=\bfi } f_\bfj \varpi_* (\phi) ([\Delta_\bfm])$ for $[\Delta_\bfm] \in \bfW_\infty$,  i.e., $\bfi\text{-Res}$ on $\bfW^*$ corresponds to  $\oplus_{ \bar\bfj=\bfi } \bfj\text{-Res}$ on $\bfW^*_\infty$ under the injection $ \varpi_*$.   
That is, we have a commutative diagram 
\begin{equation}  \label{CD:WWp}
\begin{tikzcd}
\bfW^* \ar[r,"\varpi_*"]\ar[d,"f_\bfi"] &\bfW^*_\infty= \bold{Hall}_{\infty,\Z} \ar[d,"\sum\limits_{\bfj:\bar \bfj=\bfi} f_\bfj"]
\\
\bfW^* \ar[r,"\varpi_*"] &\bfW^*_\infty= \bold{Hall}_{\infty,\Z}
\end{tikzcd} 
\end{equation}
where we identify $ \bfW^*_\infty =\bold{Hall}_{\infty,\Z}$ by what we have proved above for $p=\infty$. 
We view $\HallZ$ as a subspace of $\bold{Hall}_{\infty,\Z}$ via the injection from $\HallZ$ to $\bold{Hall}_{\infty,\Z}$ (i.e., from $\U_{p,\Z}^-$ to $\U_{\infty,\Z}^-$) which sends $f_\bfi$ to $\sum_{\bar \bfj=\bfi} f_\bfj$. Applying the generators $f_\bfi$ repeatedly to the cyclic vector $1\in \bfW^*$ gives us $\HallZ\subset \varpi_* ( \bfW^*)$. By comparing the graded ranks of $\HallZ$ and $ \bfW^*$ (for each $n$th component) we see that $\HallZ = \varpi_*(\bfW^*)$. Plugging this back into the diagram \eqref{CD:WWp} proves \eqref{equ:isomodule} for $p>0$, whence (1).

(2) Denote by $\natural^*$ the dual of the  isomorphism  $\natural$ in \eqref{equ:Moritiso}. By Theorem \ref{thm:Hall=Schursimple} and Part (1), we have $\natural^* \imath_p (f_\bfm) = \natural^*([M_\bfm]^*) =[\Delta_\bfm]^* = \jmath_p(f_\bfm)$, i.e., the
following commutative diagram holds:
\begin{equation*}  
\begin{tikzcd}
\HallZ\ar[r,"\imath_p"]\ar[dr,"\jmath_p"] &\K^*\ar[d,"\natural^*"]
\\
 &\bfW^*
\end{tikzcd} 
\end{equation*}
Then Part (2) follows from Theorems~ \ref{thm:Hall=Schursimple} and \ref{thm:ind=std}.
\end{proof}